\documentclass[11pt,a4paper]{article}
\usepackage[utf8]{inputenc}
\usepackage[spanish,es-nodecimaldot]{babel}
\usepackage{amsmath,amssymb,amsfonts}
\usepackage{geometry}
\usepackage{enumitem}
\usepackage{xcolor}
\usepackage{tcolorbox}

\geometry{
    a4paper,
    left=20mm,
    right=20mm,
    top=25mm,
    bottom=25mm
}

\title{\textbf{FAMAF -- UNC \\ Curso de Nivelación Intensivo \\ Resolución Explicativa: Ejercicios 3 y 4}}
\author{Simulacro de Primer Parcial}
\date{Agosto 2026}

\begin{document}

\maketitle

\hrule
\vspace{0.5cm}

\section*{Ejercicio 3}

\begin{tcolorbox}[colback=gray!5!white,colframe=gray!75!black,title=\textbf{Enunciado}]
En un concierto benéfico en un club se vendieron todas las entradas y se recaudaron \$230000. Los precios de las entradas eran de \$500 para menores de 10 años, y de \$3000 para adultos. La capacidad del establecimiento era de 160 personas.
\begin{enumerate}[label=\textbf{\alph*)}]
    \item Plantear un sistema de ecuaciones que describa el problema, indicando qué representa cada variable en el sistema.
    \item Calcular la cantidad de menores de 10 años y adultos que asistieron al concierto.
\end{enumerate}
\end{tcolorbox}

\subsection*{Resolución explicativa}

\subsubsection*{a) Definición de variables y planteo del sistema}

Para modelar matemáticamente la situación, comenzamos asignando variables a las incógnitas que nos pide el problema:
\begin{itemize}
    \item Sea $m$ la cantidad de entradas vendidas a menores de 10 años (número de menores asistentes).
    \item Sea $a$ la cantidad de entradas vendidas a adultos (número de adultos asistentes).
\end{itemize}

A partir de los datos del enunciado, podemos establecer dos relaciones:
\begin{enumerate}
    \item \textbf{Capacidad total:} Se vendieron todas las entradas para una capacidad de $160$ personas, por lo que la suma de asistentes debe ser:
    \begin{equation}
        m + a = 160
    \end{equation}
    \item \textbf{Recaudación total:} Cada menor abonó \$500 y cada adulto \$3000, acumulando un total de \$230000:
    \begin{equation}
        500m + 3000a = 230000
    \end{equation}
\end{enumerate}

Por lo tanto, el sistema de ecuaciones lineales de $2 \times 2$ que describe el problema es:
\begin{equation*}
\begin{cases}
    m + a = 160 \\
    500m + 3000a = 230000
\end{cases}
\end{equation*}

\subsubsection*{b) Resolución del sistema}

Podemos simplificar la ecuación de recaudación dividiendo ambos miembros por $500$:
\begin{align*}
    \frac{500m + 3000a}{500} &= \frac{230000}{500} \\
    m + 6a &= 460
\end{align*}

Ahora el sistema equivalente resulta:
\begin{equation*}
\begin{cases}
    m + a = 160 \quad &\text{(1)} \\
    m + 6a = 460 \quad &\text{(2)}
\end{cases}
\end{equation*}

Despejamos $m$ de la ecuación (1):
\begin{equation*}
    m = 160 - a
\end{equation*}

Sustituimos la expresión de $m$ en la ecuación (2):
\begin{align*}
    (160 - a) + 6a &= 460 \\
    160 + 5a &= 460 \\
    5a &= 460 - 160 \\
    5a &= 300 \\
    a &= \frac{300}{5} \\
    a &= 60
\end{align*}

Con el valor de $a = 60$, calculamos la cantidad de menores $m$:
\begin{equation*}
    m = 160 - 60 = 100
\end{equation*}

\noindent \textbf{Conclusión:} Asistieron \textbf{100 menores de 10 años} y \textbf{60 adultos} al concierto.

\vspace{0.8cm}
\hrule
\vspace{0.8cm}

\section*{Ejercicio 4}

\begin{tcolorbox}[colback=gray!5!white,colframe=gray!75!black,title=\textbf{Enunciado}]
Considerar la siguiente ecuación fraccionaria:
\begin{equation*}
    \frac{x^{3}+2x^{2}-4x-8}{x^{4}-5x^{2}+4} = -2
\end{equation*}
\begin{enumerate}[label=\textbf{\alph*)}]
    \item ¿En qué valores de $x$ no está definida la ecuación? ¿Por qué?
    \item Factorizar, simplificar la ecuación, y encontrar todas sus soluciones.
\end{enumerate}
\end{tcolorbox}

\subsection*{Resolución explicativa}

\subsubsection*{a) Dominio y valores donde no está definida la ecuación}

Una expresión algebraica fraccionaria no está definida para aquellos valores de la variable que anulen el denominador, puesto que la división por cero no está permitida en los números reales.

Buscamos los valores que hacen cero el denominador:
\begin{equation*}
    x^4 - 5x^2 + 4 = 0
\end{equation*}

Esta es una ecuación bicuadrática. Realizamos el cambio de variable $u = x^2$ (con $u \ge 0$):
\begin{equation*}
    u^2 - 5u + 4 = 0
\end{equation*}

Factorizamos el trinomio buscando dos números que multiplicados den $4$ y sumados den $-5$:
\begin{equation*}
    (u - 4)(u - 1) = 0
\end{equation*}

Por lo tanto, las soluciones para $u$ son $u_1 = 4$ y $u_2 = 1$. Volviendo a la variable original $x$:
\begin{itemize}
    \item Si $x^2 = 4 \implies x = 2 \quad \text{o} \quad x = -2$
    \item Si $x^2 = 1 \implies x = 1 \quad \text{o} \quad x = -1$
\end{itemize}

\noindent \textbf{Conclusión:} La ecuación \textbf{no está definida} para $x \in \{-2, -1, 1, 2\}$, ya que para estos valores el denominador se anula provocando una indeterminación.

\subsubsection*{b) Factorización, simplificación y conjunto solución}

Primero factorizamos de manera completa el numerador y el denominador:

\begin{itemize}
    \item \textbf{Factorización del numerador:} $P(x) = x^3 + 2x^2 - 4x - 8$. \\
    Aplicamos factor común por grupos:
    \begin{align*}
        P(x) &= x^2(x + 2) - 4(x + 2) \\
             &= (x^2 - 4)(x + 2) \\
             &= (x - 2)(x + 2)(x + 2) \\
             &= (x - 2)(x + 2)^2
    \end{align*}

    \item \textbf{Factorización del denominador:} $Q(x) = x^4 - 5x^2 + 4$. \\
    Usando las raíces obtenidas en el inciso anterior:
    \begin{align*}
        Q(x) &= (x^2 - 4)(x^2 - 1) \\
             &= (x - 2)(x + 2)(x - 1)(x + 1)
    \end{align*}
\end{itemize}

Reemplazamos las expresiones factorizadas en la ecuación para $x \notin \{-2, -1, 1, 2\}$:
\begin{equation*}
    \frac{(x - 2)(x + 2)^2}{(x - 2)(x + 2)(x - 1)(x + 1)} = -2
\end{equation*}

Simplificamos cancelando los factores comunes $(x - 2)$ y $(x + 2)$:
\begin{align*}
    \frac{x + 2}{(x - 1)(x + 1)} &= -2 \\
    \frac{x + 2}{x^2 - 1} &= -2
\end{align*}

Multiplicamos ambos miembros por $(x^2 - 1)$ (sabiendo que $x \neq \pm 1$):
\begin{align*}
    x + 2 &= -2(x^2 - 1) \\
    x + 2 &= -2x^2 + 2
\end{align*}

Igualamos a cero pasando todos los términos al primer miembro:
\begin{align*}
    2x^2 + x + 2 - 2 &= 0 \\
    2x^2 + x &= 0
\end{align*}

Sacamos factor común $x$:
\begin{equation*}
    x(2x + 1) = 0
\end{equation*}

Esto genera dos posibles soluciones:
\begin{enumerate}
    \item $x_1 = 0$
    \item $2x + 1 = 0 \implies 2x = -1 \implies x_2 = -\frac{1}{2}$
\end{enumerate}

\noindent \textbf{Verificación de restricciones:} 
Comprobamos que ni $x = 0$ ni $x = -\frac{1}{2}$ pertenezcan al conjunto de valores no permitidos $\{-2, -1, 1, 2\}$. Como ambos valores están dentro del dominio de definición, los dos son soluciones válidas.

\vspace{0.3cm}
\noindent \textbf{Conjunto solución:}
\begin{equation*}
    S = \left\{ -\frac{1}{2}, \, 0 \right\}
\end{equation*}

\end{document}
